\section{Normalized angular sectors and optical selection}
\label{sec:angular-selection}

Let $R=K^{-1/2}$, $u=\chi/R$, and
$S_K(\chi)=R\sin u$.  The spatial metric is
$\gamma=d\chi^2+S_K^2d\Omega^2$.
Write the unit-sphere embedding coordinates as
$X^0=\cos u$, $X^i=\sin u\,n^i$.
For a real symmetric tracefree $4\times4$ matrix $A$, define
\begin{equation}
 Q_A=\mathcal N X^I A_{IJ}X^J,
 \qquad \mathcal N=\sqrt{\frac{6}{\pi^2R^3}}.
 \label{eq:normalized-harmonic}
\end{equation}
The ambient polynomial is homogeneous harmonic of degree two; hence
$D^2Q_A=-8KQ_A$.  The fourth moment
$\langle X_I X_J X_K X_L\rangle=
(\delta_{IJ}\delta_{KL}+\delta_{IK}\delta_{JL}
+\delta_{IL}\delta_{JK})/24$ gives
\begin{equation}
 \int_{S^3(R)}Q_AQ_B\,dV_\gamma=\operatorname{Tr}(AB).
 \label{eq:harmonic-normalization}
\end{equation}
A Frobenius-orthonormal matrix basis therefore supplies all nine normalized
$n=2$ harmonics.  Decomposing
$A_{ij}=C_{ij}-A_{00}\delta_{ij}/3$, $\operatorname{Tr}C=0$, yields
\begin{align}
 Q_A/\mathcal N={}&A_{00}(\cos^2u-\tfrac13\sin^2u)
 +2A_{0i}\sin u\cos u\,n^i\nonumber\\
 &+\sin^2u\,n^iC_{ij}n^j.
 \label{eq:angular-decomposition}
\end{align}
These are the $1+3+5$ dimensions of the observer-centered angular sectors
$\ell=0,1,2$.  The spatial harmonic index $n$ on $S^3$ and the angular
index $\ell$ on the observer's sky are distinct.

For a screen $E$ satisfying $E^TE=I_2$, $E^Tn=0$, the tracefree spatial
Hessian of $f(\chi)Y_{\ell m}$ is
\begin{equation}
 (D_{\hat A}D_{\hat B}[fY])_{\mathrm{TF}}
 =\frac{f}{S_K^2}
 \left(\nabla_A\nabla_BY-\frac12s_{AB}\Delta_\Omega Y\right).
 \label{eq:screen-hessian}
\end{equation}
The radial connection contribution is proportional to $I_2$.
For $\ell=1$, $\nabla_A\nabla_BY=-s_{AB}Y$ and
$\Delta_\Omega Y=-2Y$, so the bracket vanishes.
The $\ell=0$ bracket also vanishes.  Thus the four monopole/dipole
spatial modes do not source first-order smooth cosmological shear on
observer-centered radial background rays.  This statement does not
remove their trace, redshift or timing responses, and does not remove
shear produced by the galaxy encounter.

For a normalized quadrupole $Q_C=\mathcal N\sin^2u\,n^TCn$ with
$\operatorname{Tr}(C^2)=1$, let
\begin{equation}
 H_C(n)=2E^TCE+(n^TCn)I_2.
 \label{eq:quadrupole-screen-shape}
\end{equation}
This matrix is symmetric and tracefree.  The exact cancellation of the
radial factor gives
\begin{equation}
 (D_{\hat A}D_{\hat B}Q_C)_{\mathrm{TF}}
 =\mathcal N K H_C(n).
 \label{eq:quadrupole-radial-cancellation}
\end{equation}

\section{State-derived tracefree propagation}
\label{sec:tracefree-transfer}

In the convention
$ds^2=a^2[-(1+2\Phi)d\eta^2+(1-2\Psi)\gamma_{ij}dx^idx^j]$,
write the current six temporal metric responses as
$\phi_j(\eta)$ and $\psi_j(\eta)$, and let
$w_j=\phi_j+\psi_j$.
The initial coordinates are
$(\zeta,\Pi_2,\delta_m^{\rm ADM},\delta_r^{\rm ADM},
v_m^{\rm ADM},v_r^{\rm ADM})$ for each normalized spatial harmonic.
No equality between the two metric potentials is imposed.

The background physical angle-to-separation map between vertices $a,b$ is
$d_{ab}=a_bS_K(\chi_b-\chi_a)$.
The physical initialization $D'_a=\omega_aI_2$, with one
observer-normalized affine parameter, becomes a unit conformal derivative
at $\chi_a$.  Linearizing the Sachs equation gives
\begin{equation}
 (\partial_\chi^2+K)\delta\widetilde D_{ab}^{\rm TF}
 =-S_K(\chi-\chi_a)
 [D_\perp D_\perp(\Phi+\Psi)]_{\rm TF}.
 \label{eq:linearized-sachs-tf}
\end{equation}
Define shear by $\delta D_{ab}^{\rm TF}=-d_{ab}\Gamma_{ab}$.
For an orthonormal STF three-matrix basis $C_A$, $A=1,\ldots,5$,
the result is
\begin{equation}
 \Gamma_{ab}(n)=\sum_{j=1}^{6}\sum_{A=1}^{5}
 g_{ab,j}y_{jA}H_{C_A}(n),
 \label{eq:shear-response}
\end{equation}
\begin{equation}
 g_{ab,j}=\frac{\mathcal NK}{S_K(\chi_b-\chi_a)}
 \int_{\chi_a}^{\chi_b}S_K(\chi_b-\chi)S_K(\chi-\chi_a)
 w_j(\eta_o-\chi)\,d\chi.
 \label{eq:shear-kernel}
\end{equation}
The three maps $ol$, $os$, and $ls$ are independently initialized.
The final scale $a_b$ cancels from this fractional shear, and no additional
redshift multiplier belongs in the kernel.

At fixed measured redshift, endpoint motion and moving-vertex terms
multiply the background identity at first order.  Corrections to a
first-order shear integral from those shifts are second order.  Local
aberration is conformal on an infinitesimal celestial sphere and supplies
no first-order spin-two term.  These facts establish the tracefree result;
they do not replace the frequency, radial-event and local-scale terms
required for a generic trace response.

The full retained state has 54 coefficients, of which at most 30 enter
this smooth-shear response.  For redshift/vertex labels $r$,
the five angular coefficients satisfy $S_{rA}=G_{rj}y_{jA}$,
$G_{rj}=g_{r,j}$.  A left-null vector of the actual kernel design
therefore supplies a shared-state relation.  Its observational relevance
depends on conditioning, units and the specified lens/source nuisance model.
If all six state components share a single spatial quadrupole $C$,
then $\Gamma_{ab}=(g_{ab}\cdot y)H_C$, and
\begin{equation}
 \gamma_{1,ab}\gamma_{2,cd}-\gamma_{2,ab}\gamma_{1,cd}=0.
 \label{eq:common-profile-alignment}
\end{equation}
This profile restriction is not automatic for a generic six-by-five
quadrupole preparation, and the deterministic identity is not a
stochastic covariance prediction.

\section{A complete conditional timing response on a nodal ray}
\label{sec:nodal-timing}

For a weak thin local potential $\widehat\psi$, use the standard
distance-matrix arrival-time function
\begin{align}
 \beta&=D_{os}^{-1}X_s,
 &T&=(1+z_l)D_{ol}^TD_{ls}^{-1}D_{os},\nonumber\\
 ct(\theta)&=\tfrac12(\theta-\beta)^TT(\theta-\beta)
 -(1+z_l)\widehat\psi(D_{ol}\theta)+f(\beta).
 \label{eq:matrix-fermat}
\end{align}
This optical construction and its clock normalization are established
distance-matrix lens theory \cite{Fleury2021,Schneider2019}.
Put $D_{ab}=d_{ab}(I_2+E_{ab})$ and
$T_0=(1+z_l)d_{ol}d_{os}/d_{ls}$.
At first order $\delta T/T_0=E_{ol}+E_{os}-E_{ls}$.
For a stationary image, fixed measured redshifts, fixed transverse
physical source $X_s$, and fixed local physical potential, the source
variation is $\delta\beta=-E_{os}\beta$.
Using the background lens equation and writing $v_i=\theta_i-\beta$,
\begin{equation}
 \frac{c\delta t_i-\delta f}{T_0}
 =\tfrac12v_i^T(E_{os}-E_{ol}-E_{ls})v_i
 +v_i^T(E_{os}-E_{ol})\beta.
 \label{eq:stationary-delay-response}
\end{equation}
The expression includes the varying impact in the local potential.
Image-position derivatives drop out by stationarity; the common $f$
cancels in relative delays.

Choose $C_\star=\operatorname{diag}(1,-1,0)/\sqrt2$, the central ray
$n=e_z$, and screen $(e_x,e_y)$.
Here $Q$ and its full first gradient vanish on the ray, the transverse
Hessian has zero trace, and
$H_{C_\star}=\sqrt2\operatorname{diag}(1,-1)$.
The first-order frequency, radial-event, velocity, local-scale and Jacobi
trace responses therefore vanish.  The first-order physical maps are
fully specified by $E_{ab}=-\Gamma_{ab}$ for this preparation.

Take a local SIS, $\widehat\psi(x)=b|x|$, and a source
$X_s=s e$, $e=(\cos\varphi,\sin\varphi)$, with fixed $b,s>0$.
The background values are $\beta=s/d_{os}$,
$\theta_E=bd_{ls}/d_{os}$, and
$\theta_\pm=(\beta\pm\theta_E)e$.
For $0<\beta<\theta_E$ and fixed noncritical branches,
\begin{equation}
 \Delta t_0=t_- -t_+
 =\frac{2(1+z_l)bs}{c}\frac{d_{ol}}{d_{os}}.
 \label{eq:sis-background-delay}
\end{equation}
Since $v_\pm=\pm\theta_E e$, the quadratic term in
Eq.~\eqref{eq:stationary-delay-response} cancels between images:
\begin{equation}
 \boxed{\frac{\delta\Delta t}{\Delta t_0}
 =e^T(\Gamma_{os}-\Gamma_{ol})e
 =\sqrt2[(g_{os}-g_{ol})\cdot y]\cos(2\varphi).}
 \label{eq:quadrupole-delay-law}
\end{equation}
The $ls$ shear cancels for this fixed-physical-source SIS experiment.
The result predicts the angular sign reversal and null directions for
the declared preparation; it does not select the physical amplitude $y$.
It is not a generic prediction for arbitrary galaxy profiles, fixed
angular sources, finite-amplitude metric completions, or higher-order
light propagation.
